\chapter{Order Gateway and Order Management}\label{ch:ll:order-gateway-and-order-management}

A strategy cancels a resting buy order of 100 lots and, in the same microsecond, sends its replacement one tick lower. On the
way to the exchange the cancel crosses a fill: a seller has just taken the old order. The exchange rejects the cancel as too
late, accepts the new order, and a second seller takes it too. For the few microseconds between the strategy's decision and the
fill report, the firm held twice the position it believed it held, and a gateway that counted only the orders it knew to be
live would have let the strategy send a third. Nothing in this story is a bug in the ordinary sense: every message was valid,
and every party did what the protocol says. It is what happens when two machines a network apart each change an order's state
without waiting for the other. This chapter builds the component that keeps the firm's view of its orders honest while
this goes on: the \omterm{def:ll:order-gateway-and-order-management:gateway}{order gateway}, with a state machine that has names for ``I asked and have not heard back'', an exposure count
that assumes the worst about every order still in flight, a throttle that keeps the venue's message budget, and a
reconciliation against the exchange's own copy of the day.

\section{The order state machine}

\begin{definition}[Order gateway, order state machine]\label{def:ll:order-gateway-and-order-management:gateway}
An \emph{order gateway}\index{order gateway} is the component between a firm's strategies and a venue's order-entry session:
it turns the strategies' requests into the venue's messages, the venue's reports into order states, and refuses what the firm's
limits or the venue's rules forbid. An \emph{order state machine}\index{order state machine} is the set of states an order can
be in, as the firm sees it, with the events that move it from one to another; the firm's view includes the \emph{pending}
states, entered when the firm has sent a request and not yet heard its answer.
\end{definition}

The venue's order has few states: live, partially filled, and then filled, cancelled or expired (Book 10's simulator, its
\texttt{PROTOCOL.md}). The firm's has more, because the firm is always a little behind. After it sends a new order, the order is
\emph{pending new} until the acknowledgement arrives; after it asks for a cancel or a replace, it is \emph{pending cancel} or
\emph{pending replace} until the answer arrives, and in the meantime the order can still fill. FIX names the same states in its
order status field (\emph{PendingNew}, \emph{PendingCancel}, \emph{PendingReplace}, next to \emph{New}, \emph{PartiallyFilled},
\emph{Filled}, \emph{Canceled} and \emph{Rejected}; chapter~15): the protocol itself admits that the two sides disagree for a
while. \Cref{fig:ll:order-gateway-and-order-management:states} draws the firm's machine.

\begin{omfigure}
\begin{tikzpicture}[font=\scriptsize,>=Stealth,st/.style={draw,rounded corners=3pt,minimum width=1.9cm,minimum height=0.62cm,
  align=center,fill=omDef!12},pend/.style={st,fill=omMeth!18,dashed},term/.style={st,fill=gray!12,double}]
\node[pend] (pn) at (0,0) {pending new};
\node[st] (live) at (3.3,0) {live};
\node[st] (part) at (6.6,0) {partial};
\node[term] (fil) at (9.9,0) {filled};
\node[term] (rej) at (0,-2.4) {rejected};
\node[pend] (pc) at (3.3,-2.4) {pending cancel};
\node[pend] (pr) at (6.6,-2.4) {pending replace};
\node[term] (cxl) at (9.9,-2.4) {cancelled};
\draw[->] (pn) -- node[above]{ack} (live);
\draw[->] (live) -- node[above]{fill} (part);
\draw[->] (part) -- node[above]{last fill} (fil);
\draw[->] (pn) -- node[left]{reject} (rej);
\draw[->] (live) -- node[left,align=right]{cancel\\ sent} (pc);
\draw[->] (part) -- node[right,align=left]{replace\\ sent} (pr);
\draw[->] (pr.north west) -- node[below,sloped]{replaced} (live.south east);
\draw[->] (pr) -- node[above]{cancelled} (cxl);
\draw[->,omThm,thick] (pr.north east) -- node[below,sloped,text=omThm]{last fill} (fil.south west);
\draw[->] (pc.south) -- ++(0,-0.7) -| node[pos=0.25,above]{cancelled} (cxl.south);
\draw[->,draw=omThm,thick] ([xshift=-4mm]pc.south) |- node[pos=0.75,below,text=omThm]{last fill, then ``too late''}
  (11.5,-3.9) |- (fil.east);
\draw[->,omThm,thick] (pc.west) .. controls +(-0.9,0.8) and +(-0.2,0.9) .. node[pos=0.3,above left,text=omThm,inner sep=1pt]{partial fill} ([xshift=-6mm]pc.north);
\end{tikzpicture}

\omcaption{The firm's view of an order. Dashed boxes are the pending states, entered when a request is sent and left when its
answer arrives; in red, the fills that reach an order after the firm has asked to cancel it. Not drawn: every open state can also be
cancelled by the venue on its own (a lost connection, a halt, an expiry), and pending new can fill or be cancelled before
its acknowledgement. Only the terminal states, doubled, are final; the table of \cref{lst:ll:order-gateway-and-order-management:table} is the whole
machine.}\label{fig:ll:order-gateway-and-order-management:states}
\end{omfigure}

The build writes the machine as a table, one row per state and one column per event, with a marker for the transitions that
are not allowed (\cref{lst:ll:order-gateway-and-order-management:table}). A table has three advantages over a nest of
conditions. It is complete by construction: every pair of state and event has an entry, so the question ``what if a fill
arrives while a replace is pending?'' has an answer that someone wrote down. It is data: the Python reference, the C++ gateway
and the Rust gateway each hold a copy, and the tests check that the three copies are the same table, entry for entry,
against \texttt{data/table.txt}. And an impossible event is detected, not absorbed: a report that the table does not allow
is returned as an error, which is how the first version of the build found that it had no arrow from \emph{live} to
\emph{cancelled}. The venue cancels orders on its own when a session is lost, when trading halts or when an order expires, and
the conformance test against Book 10's live server, which cuts a connection on purpose, hit the missing entry on its first run.

\omcode{code/firm/ordergw/cpp/firm_ordergw.hpp}{28}{38}{The order state machine as a table: the next state for every state
and event, and a marker for what is not allowed.}\label{lst:ll:order-gateway-and-order-management:table}

A replace deserves a word. The simulator, like most venues, gives the replaced order a new client identifier, and the gateway
keeps the order's history under it; while the replace is pending, the order can fill under its old identifier, so the gateway
keeps both until the answer arrives. Whether the order keeps its place in the queue is the venue's rule (the report says
whether priority was kept); the gateway only records it.

\section{Acknowledgements, in-flight orders and races}

\begin{definition}[In-flight order, cancel--fill race]\label{def:ll:order-gateway-and-order-management:inflight}
An \emph{in-flight order}\index{in-flight order} is an order in a pending state: a request about it has been sent and its
answer has not arrived, so the order may have changed at the venue in ways the firm does not know yet. A \emph{cancel--fill
race}\index{cancel--fill race} occurs when a fill and a request to cancel or replace the same order cross: the fill happens at
the venue before the request arrives there, and its report reaches the firm after the request was sent.
\end{definition}

\Cref{fig:ll:order-gateway-and-order-management:race} draws one. The strategy sends its cancel at time $c$; it arrives at the
venue a one-way latency $L$ later; each fill report takes $d$ to come back. Any fill in the interval from $c - d$ to $c + L$
races the cancel: the fills before $c - d$ were known when the strategy decided, and the ones after $c + L$ cannot happen,
because the order is gone. If fills arrive at a Poisson rate $\lambda$ on a resting order, the probability that a cancel meets
at least one is
\[
P(\text{race}) = 1 - e^{-\lambda (L + d)} \approx \lambda (L + d),
\]
as long as the order has rested longer than $L + d$ before the cancel (an order cancelled sooner has had less time exposed, and
the weekend problem computes the exact probability). The race is not a failure of the venue or of the gateway. It is the
latency of the path, seen from the order's point of view, and the only ways to make it rarer are to shorten the path or to
cancel less.

\begin{omfigure}
\begin{tikzpicture}[font=\scriptsize,>=Stealth]
\draw[thick] (0,0) node[above]{strategy and gateway} -- (0,-4.4);
\draw[thick] (6.4,0) node[above]{venue} -- (6.4,-4.4);
\fill[omThm!12] (6.3,-1.1) rectangle (6.5,-2.9);
\draw[gray] (6.3,-1.1) -- (6.5,-1.1) node[pos=0,left,gray]{$c - d$};
\draw[->,omDef,thick] (0,-1.9) node[left]{cancel sent, $c$} -- node[above,sloped,pos=0.22]{cancel, latency $L$} (6.4,-2.9);
\draw[->,omThm,thick] (6.4,-1.6) node[right]{fill} -- node[above,sloped,pos=0.4]{fill report, delay $d$} (0,-2.9);
\node[left,align=right] at (0,-2.9) {fill arrives:\\ state pending cancel};
\draw[->,omDef] (6.4,-2.9) -- node[below,sloped,pos=0.35]{``too late''} (0,-3.9);
\node[right,align=left,omDef] at (6.4,-3.25) {cancel arrives at $c + L$:\\ nothing left to cancel};
\node[left] at (0,-3.9) {filled};
\draw[decorate,decoration={brace,amplitude=4pt}] (7.1,-1.1) -- node[right=4pt,align=left]{a fill here races the\\ cancel: from $c - d$\\ to $c + L$} (7.1,-2.9);
\end{tikzpicture}

\omcaption{A \omterm{def:ll:order-gateway-and-order-management:inflight}{cancel--fill race}. The cancel leaves at $c$ and arrives $L$ later; a fill at the venue between $c - d$ and $c + L$
is reported after the cancel was sent. The gateway, in the state \emph{pending cancel}, applies the fill and then the ``too
late'' answer, and ends in \emph{filled}.}\label{fig:ll:order-gateway-and-order-management:race}
\end{omfigure}

The build measures it. \texttt{firm\_ordergw\_sim.py} runs the Python gateway against Book 10's matching engine: a strategy
keeps one buy order working, holds it for an exponentially distributed time of mean \qty{1}{\milli\second} and cancels it; a
counterparty sells into it at random, 200 times a second; every report takes \qty{20}{\micro\second} to come back.
\Cref{fig:ll:order-gateway-and-order-management:races} gives the fraction of cancels that met a fill, for one-way latencies
from \qty{10}{\micro\second} to \qty{3}{\milli\second}. At \qty{100}{\micro\second}, 2.1\% of the cancels raced, against 2.2\%
from the exact formula; at \qty{10}{\micro\second} it is 0.65\%. At \qty{1}{\milli\second} the simple formula would say 18\%
and the engine says 12\%: the orders are held for about as long as the path takes, so the window is often shorter than $L + d$.

\begin{omfigure}
\begin{tikzpicture}
\begin{axis}[width=11cm,height=6cm,xmode=log,xmin=8,xmax=4000,ymin=0,ymax=0.5,
  xlabel={one-way latency of the cancel $L$ ($\mu$s)},ylabel={fraction of cancels that raced},
  tick label style={font=\small},label style={font=\small},grid=major,grid style={gray!25},
  legend pos=north west,legend style={font=\scriptsize,cells={anchor=west}},table/col sep=comma,
  yticklabel style={/pgf/number format/fixed}]
\addplot[omMeth,thick,dashed] table[x=latency_us,y=rested]{figdata/low-latency/21-order-gateway-and-order-management/race_curves.csv};
\addlegendentry{$1 - e^{-\lambda (L + d)}$ (orders rested at least $L + d$)}
\addplot[omDef,thick] table[x=latency_us,y=exact]{figdata/low-latency/21-order-gateway-and-order-management/race_curves.csv};
\addlegendentry{exact, orders held \qty{1}{\milli\second} on average}
\addplot[only marks,mark=*,omThm,mark size=2.2pt] table[x=latency_us,y=simulated]{figdata/low-latency/21-order-gateway-and-order-management/races.csv};
\addlegendentry{gateway against the matching engine}
\end{axis}
\end{tikzpicture}

\omcaption{\omterm{def:ll:order-gateway-and-order-management:inflight}{Cancel--fill races} against the latency of the cancel, with fills at $\lambda = 200$ per second and reports
delayed $d = \qty{20}{\micro\second}$. The points are the Python gateway against Book 10's matching engine (two seeded runs of
2\,000 orders each per point); the solid curve is the exact probability for orders held \qty{1}{\milli\second} on average,
the dashed one the formula for orders that have rested at least $L + d$ when cancelled. Data: \texttt{fig\_races.py}
(deterministic).}\label{fig:ll:order-gateway-and-order-management:races}
\end{omfigure}

What the gateway does with a race decides whether it is harmless. The hook's story is the classic failure: the gateway counted
an order as gone when it \emph{sent} the cancel, so the replacement fitted under the position limit, and both filled. The rule
that prevents it is \emph{worst-case accounting}: until the venue has answered, an order counts for everything it could still
do. An order pending cancel counts its whole remaining quantity; an order pending replace counts the larger of its old and its
new quantity, because either may fill; an order pending new counts in full. The long exposure is then the position plus the
worst remaining quantity of every open buy order, and a new order is sent only if it fits under the limit on top of it. With a
limit of 100 lots, the hook's gateway refuses the replacement until the cancel's answer arrives, and the firm's position never
exceeds 100. The price is paid in opportunity: for one round trip, the strategy cannot use the quantity it is trying to
withdraw.

The gateway keeps this exposure as running sums, one per side, updated whenever an order changes: take the order's worst
remaining quantity out, change the order, put it back (\cref{lst:ll:order-gateway-and-order-management:change}). The check
before a new order is then one addition and one comparison (\cref{lst:ll:order-gateway-and-order-management:new}), whatever
the number of open orders. A gateway that recomputes the exposure by walking its orders pays for every one of them on every
new order: \cref{fig:ll:order-gateway-and-order-management:exposure} measures about \qty{0.4}{\micro\second} for a thousand
open orders kept in an array and about \qty{1.6}{\micro\second} in a hash map, where the running sums cost less than the timer
can resolve. A market maker quoting a few hundred instruments has thousands of open orders all day.

\omcode{code/firm/ordergw/cpp/firm_ordergw.hpp}{218}{225}{Worst-case exposure kept as running sums: every change to an order
goes through \texttt{change}.}\label{lst:ll:order-gateway-and-order-management:change}

\omcode{code/firm/ordergw/cpp/firm_ordergw.hpp}{83}{98}{A new order: the exposure check against the worst case, the throttle,
then the venue's message written in place with the codec of chapter~16.}\label{lst:ll:order-gateway-and-order-management:new}

\begin{omfigure}
\begin{tikzpicture}
\begin{axis}[width=10cm,height=5.8cm,xmode=log,ymode=log,xmin=0.7,xmax=15000,ymin=0.5,ymax=40000,
  xlabel={open orders},ylabel={exposure check (ns)},tick label style={font=\small},label style={font=\small},
  grid=major,grid style={gray!25},legend pos=north west,legend style={font=\scriptsize,cells={anchor=west}},
  table/col sep=comma]
\addplot[omThm,thick,mark=square*] table[x=open_orders,y=map_ns]{figdata/low-latency/21-order-gateway-and-order-management/measured_exposure.csv};
\addlegendentry{scan of a hash map}
\addplot[omMeth,thick,mark=*] table[x=open_orders,y=array_ns]{figdata/low-latency/21-order-gateway-and-order-management/measured_exposure.csv};
\addlegendentry{scan of an array}
\end{axis}
\end{tikzpicture}

\omcaption{The cost of the worst-case exposure check before each new order, recomputed by scanning the open orders, against
their number (median of timed calls). Kept as running sums, as in \texttt{firm.ordergw}, it costs one addition, below the
timer's resolution at every size. Measured on a laptop (Intel Core Ultra 7 155H) under WSL2. Data:
\texttt{bench\_gateway.py}.}\label{fig:ll:order-gateway-and-order-management:exposure}
\end{omfigure}

The same benchmark times the gateway's whole work for each step of an order's life: writing a new order, applying its
acknowledgement, a partial fill, writing the cancel and applying its answer, each decoded from the venue's binary report.
Every step costs 4 to \qty{5}{\nano\second} at the median and 5 to \qty{6}{\nano\second} at the 99th
percentile: the gateway's logic is not where the microseconds of an order's path go. They go to the network and the kernel
(chapter~13 and One Quant Book~14), and to the risk checks of the next chapter.

\section{Throttles and message budgets}

\begin{definition}[Token bucket]\label{def:ll:order-gateway-and-order-management:bucket}
A \emph{token bucket}\index{token bucket} is a rate limiter that holds up to $b$ tokens and gains $r$ tokens a second; a
message may be sent only if a token is available, and sending it spends one. It allows bursts of up to $b$ messages and, over
any interval of length $T$, at most $b + rT$ messages.
\end{definition}

Venues limit how many messages a session may send: a rate per second, sometimes a burst, sometimes a weight per message type
or a ratio of orders to trades (the message throttle of One Quant Book~11, chapter~10, and Book 10's simulator, which rejects
the excess with reason \texttt{T}). A gateway that relies on the venue to enforce the limit learns about it from rejections,
after the fact, and a venue may disconnect a session that keeps exceeding it. The gateway therefore keeps its own budget, set
a little below the venue's, and refuses locally what would be rejected remotely. The build's bucket counts in integer
nano-tokens, so that the Python, C++ and Rust gateways agree to the last message on the same journal
(\cref{lst:ll:order-gateway-and-order-management:bucket}).

\omcode{code/firm/ordergw/cpp/firm_ordergw.hpp}{43}{50}{The token bucket in integer nano-tokens: refill for the time elapsed,
then spend one token or refuse.}\label{lst:ll:order-gateway-and-order-management:bucket}

What to do with a refused message is the strategy's decision, not the gateway's. A cancel refused by the throttle is the worst
case: the order stays in the market while the strategy wanted it out. Gateways therefore keep part of the budget for cancels,
or let cancels through ahead of new orders when the budget is short; the build's gateway reports the reason for each refusal
(throttle, exposure, state) so that the strategy can tell a full budget from a limit. The build's replace fixture runs a
strategy that replaces its order every millisecond on average against a budget of 300 messages a second with a burst of two:
the gateway refuses 111 of its requests for the budget, and the strategy retries.

\section{Sessions: login, sequence numbers, cancel on disconnect}

The venue's order-entry session (Book 10, chapter~26) carries the gateway's messages. Book 10's simulator speaks SoupBinTCP:
the client logs in with a user name and a password; its messages go unsequenced; the server's reports are sequenced, numbered
from one per session, but the numbers are not on the wire: each side counts. A client that loses its connection logs in again
asking for the next report it has not seen, and the server replays every report from that number. The client's half of the
protocol is short (\cref{lst:ll:order-gateway-and-order-management:session}): count the sequenced messages, ask for the next
one at login, check that the server starts there, send a \omterm{def:ll:protocols-i-fix:session}{heartbeat} after a second of silence, and consider the server dead
after fifteen.

\omcode{code/firm/ordergw/firm_ordergw.py}{291}{305}{The client side of the order-entry session: a login accepted at the
wrong sequence number is an error, not a detail.}\label{lst:ll:order-gateway-and-order-management:session}

A lost connection is where worst-case accounting earns its keep. While the gateway is disconnected, it does not know what its
orders are doing: they may be filling, or the venue may have cancelled them. Most venues offer, and many require, a cancel on
disconnect: the venue cancels the session's resting orders when it detects that the connection is lost (the simulator does it
by default, with cancel reason \texttt{D}). The gateway still counts them in full until the reports say otherwise, because
cancel on disconnect is a feature of the venue, not a guarantee to the firm: it acts only on the orders it covers, when it
detects the loss, and a fill can precede it. The conformance test does exactly this against Book 10's live server on the
loopback interface: three orders go out, one fills, the connection is cut, the exposure stays at 300 lots while the gateway is
away, and after the new login the replayed reports cancel the other two and bring it back to 100.

\begin{dated}{2026-09}{Cancel on disconnect and drop copy at one exchange}\label{dat:ll:order-gateway-and-order-management:cme}
CME Group's client documentation (consulted September 2026) describes Cancel On Disconnect for its iLink order-entry
sessions: when a lost connection is detected, it cancels all resting futures and options orders of the disconnected session,
except good-till-cancelled and good-till-date orders, and new sessions are created with it enabled by default. Its Drop Copy
service sends real-time copies of the execution reports and acknowledgements sent over iLink sessions on a separate,
dedicated path.
\end{dated}

\section{Drop copies and reconciliation}

\begin{definition}[Drop copy]\label{def:ll:order-gateway-and-order-management:dropcopy}
A \emph{drop copy}\index{drop copy} is a separate session on which the venue sends copies of the reports of a firm's trading
sessions (at least its executions and cancellations) as it sends the originals; it is read-only, and serves as an independent
record of the firm's orders and fills.
\end{definition}

The gateway's view of the firm's orders is built from the reports it received, applied in order; any lost, duplicated or
misapplied report makes it wrong, and a wrong view is found only by comparing it with another. The \omterm{def:ll:order-gateway-and-order-management:dropcopy}{drop copy} is that other
view, from the venue, on a different connection, often read by a different process (Book 10's simulator offers one session of
this kind per firm). Reconciliation compares the two: for every order, the filled quantity and whether it was cancelled
(\cref{lst:ll:order-gateway-and-order-management:reconcile}). Each difference is a \emph{break}, and a break is investigated,
not fixed by overwriting one side: an order the \omterm{def:ll:order-gateway-and-order-management:dropcopy}{drop copy} shows filled and the gateway shows live is a lost report and an
unknown position; the opposite is a report applied twice.

\omcode{code/firm/ordergw/firm_ordergw.py}{189}{203}{Reconciliation against the drop copy: filled quantities and cancellations,
order by order; every difference is a break.}\label{lst:ll:order-gateway-and-order-management:reconcile}

In the build, every simulated run and the live test end with zero breaks, and the property tests make sure that this is not
luck: over random interleavings of fills, cancels, replaces and latencies from \qty{5}{\micro\second} to \qty{2}{\milli\second},
no order ever has more filled than ordered, no remaining quantity is negative, the position equals the sum of the fills, and
the worst-case exposure is never below the position; and with a limit of 100 lots, the position against the matching engine
never exceeds 100. The firm reconciles at several speeds: continuously against the \omterm{def:ll:order-gateway-and-order-management:dropcopy}{drop copy}, every morning against
the broker's records (One Quant Book~1, chapter~7), and the order management system above the gateway (Book~10,
chapter~20) against the portfolio's books.

\section{Tutorial: races, throttles and a reconciliation}\label{tut:ll:order-gateway-and-order-management:tut}

\begin{tutorial}
\textbf{Goal.} Run the gateway against the matching engine, count its races, replay its journal in three languages, and
reconcile it against the \omterm{def:ll:order-gateway-and-order-management:dropcopy}{drop copy} of a live session. \textbf{End state:}
\cref{fig:ll:order-gateway-and-order-management:races}, the replays equal to \texttt{data/expected.txt}, and a reconciliation
with zero breaks.
\begin{tutsteps}
\item \textbf{Races.} \texttt{python fig\_races.py} runs \texttt{firm\_ordergw\_sim.run} at six latencies and writes the
fractions of \cref{fig:ll:order-gateway-and-order-management:races}, next to the exact formula of \texttt{ll\_gateway.py}.
\item \textbf{Journals.} \texttt{make\_ordergw\_fixtures.py} records two runs (cancel and new; cancel--replace under a tight
budget) as journals of requests and reports; the C++ and Rust gateways replay them, re-making every request with the same
result, and print the summaries the Python reference printed. The C++ test also feeds the reports as encoded binary messages,
and counts no allocation.
\item \textbf{Session.} \texttt{tests/test\_live\_session.py} starts Book 10's live server, logs in a trading session, a \omterm{def:ll:order-gateway-and-order-management:dropcopy}{drop
copy} and a counterparty, cuts the trading session's connection, logs in again, and reconciles.
\item \textbf{Costs.} \texttt{python bench\_gateway.py} times each step of an order's life and the exposure check against the
number of open orders.
\end{tutsteps}
\textbf{What to change next.} Lower the replace fixture's budget to 100 messages a second and watch the refusals; make the
counterparty's fills arrive in bursts instead of at a constant rate and compare the race count with the formula.
\end{tutorial}

\section{Build: the order gateway}\label{bld:ll:order-gateway-and-order-management:ordergw}

\begin{build}
\bfield{Purpose} The firm's single path to the venue's order entry: every order from the strategy engine (chapter~20) passes
through it, after the risk gate (chapter~22); its journal is logged (chapter~23), and its state survives a failover
(chapter~24).
\bfield{Interface} Python reference \texttt{firm\_ordergw}: \texttt{TABLE}, \texttt{transition}, \texttt{TokenBucket},
\texttt{Gateway(rate, burst, max\_long, max\_short)} with \texttt{new}, \texttt{cancel}, \texttt{replace}, \texttt{on\_report},
\texttt{worst\_long}, \texttt{worst\_short}, \texttt{races}; \texttt{Session}, \texttt{frames}, \texttt{reconcile},
\texttt{replay}, \texttt{summary}; \texttt{firm\_ordergw\_sim.run} against Book 10's engine. C++20 \texttt{firm::ogw}:
\texttt{Gateway(Limits)} with \texttt{new\_order}, \texttt{cancel}, \texttt{replace} writing the venue's messages in place,
\texttt{on\_report} and \texttt{on\_wire}, \texttt{summary}. Rust \texttt{firm\_ordergw}: \texttt{Gateway}, \texttt{replay}.
\bfield{Rules} Every open order counts for the worst it can still do; nothing is sent past the exposure limit or the budget;
a report the table does not allow is an error; an order is kept until a terminal report; the three gateways hold the same
table; no allocation per message.
\bfield{Acceptance tests} \texttt{code/firm/ordergw/}: the table identical in three languages; both journals replayed to the
committed summaries in Python, C++ (from fields and from the wire) and Rust; property tests over random runs against the
matching engine; the exposure limit held against the engine; the naive gateway's double position reproduced and refused by
worst-case accounting; the race fraction within a point of the formula; the live session with cancel on disconnect, replay
and a zero-break reconciliation.
\bfield{Stretch} Mass cancel and mass quotes; a budget reserved for cancels; the session layer in C++ on Book 10's
\texttt{exchsim\_client.hpp}; \omterm{def:ll:protocols-i-fix:session}{FIX sessions} through the engine of chapter~15.
\end{build}

\begin{omsources}
\item CME Group Client Systems Wiki, \emph{Cancel on Disconnect}, and \emph{Drop Copy 4.0 Service for iLink}.
\item FIX Trading Community, FIX Latest, code set \emph{OrdStatus}.
\end{omsources}

\section{Exercises}

\begin{exercise}[$\star$]\label{exo:ll:order-gateway-and-order-management:1}
An order is pending cancel. A fill for its whole remaining quantity arrives, then the venue's rejection of the cancel as too
late. Give the order's state after each report, and say what the gateway should count as its remaining quantity in between.
\end{exercise}

\begin{exercise}[$\star$]\label{exo:ll:order-gateway-and-order-management:2}
A \omterm{def:ll:order-gateway-and-order-management:bucket}{token bucket} gains 300 tokens a second and holds at most 2; it is full at time 0. Messages are sent at 0, 0, 0, 1, 4 and
5 milliseconds. Which go through?
\end{exercise}

\begin{exercise}[$\star$]\label{exo:ll:order-gateway-and-order-management:3}
The long limit is 500 lots. The position is 200 lots long, and two buy orders of 100 lots are live. The strategy asks to
replace one of them with a buy of 300. Does the gateway send the replace? What is the largest new quantity it would send?
\end{exercise}

\begin{exercise}[$\star\star$]\label{exo:ll:order-gateway-and-order-management:4}
Fills arrive at 50 per second on a resting order, the cancel takes \qty{300}{\micro\second} to reach the venue and reports
take \qty{50}{\micro\second} to come back. What fraction of cancels race a fill? How many races should a strategy that cancels
20\,000 orders a day expect?
\end{exercise}

\begin{exercise}[$\star\star$]\label{exo:ll:order-gateway-and-order-management:5}
The \omterm{def:ll:order-gateway-and-order-management:dropcopy}{drop copy} shows two executions of 60 lots for order 17, both with match number 88\,412; the gateway shows order 17 filled
for 60. Which side is wrong, and what should the reconciliation have done?
\end{exercise}

\begin{exercise}[$\star\star$]\label{exo:ll:order-gateway-and-order-management:6}
A gateway recomputes the exposure by scanning a hash map of its open orders before each new order. With a thousand open orders,
what does that add to each order according to \cref{fig:ll:order-gateway-and-order-management:exposure}, and how does it compare
with the gateway's own work per message?
\end{exercise}

\begin{exercise}[$\star\star\star$]\label{exo:ll:order-gateway-and-order-management:7}
\emph{Coding.} Add mass cancel to the gateway: one request that cancels every open order of a side. Which states does each order
enter, what does the exposure count until the answers arrive, and how do you test it against Book 10's engine?
\end{exercise}

\begin{exercise}[$\star\star\star$]\label{exo:ll:order-gateway-and-order-management:8}
\emph{Find the flaw.} ``To keep the order table small, our gateway frees an order's slot as soon as it sends the cancel: the
strategy does not want the order any more.''
\end{exercise}

\section{Problem: The Cancel That Crossed a Fill}

\begin{problem}[{Weekend problem --- the window in which a cancel can lose}]\label{pb:ll:order-gateway-and-order-management:1}
A market maker quotes buy orders of 100 lots, holds each for a millisecond on average and then cancels it. Fills arrive at
$\lambda = 200$ per second on a resting order, reports take $d = \qty{20}{\micro\second}$ to come back, and the cancel takes
$L$ to reach the venue. Use \cref{fig:ll:order-gateway-and-order-management:races} (\texttt{races.csv}).

\medskip\noindent\textbf{Part I --- The window.}
\begin{enumerate}
\item Explain why a fill at the venue between $c - d$ and $c + L$ races a cancel sent at $c$, and why fills outside that
interval do not.
\item Derive $P(\text{race}) = 1 - e^{-\lambda(L + d)}$ for an order that has rested at least $L + d$ when the cancel is sent.
\item Compute it for $L = \qty{10}{\micro\second}$, \qty{100}{\micro\second} and \qty{1}{\milli\second}.
\item Show that it is approximately $\lambda (L + d)$ when that product is small, and say how small.
\end{enumerate}

\medskip\noindent\textbf{Part II --- The measurement.}
\begin{enumerate}[resume]
\item Read the simulated fractions at the same three latencies from the figure's data.
\item Why does the simple formula overestimate them at \qty{1}{\milli\second}?
\item For an order held a time $h$, the window is $\min(h, L + d)$. What happens to the race probability as $L$ grows much
larger than the holding time?
\item The committed cancel journal (latency \qty{100}{\micro\second}) has 6 races in 245 cancels. Is that consistent with the
figure?
\end{enumerate}

\medskip\noindent\textbf{Part III --- The damage.}
\begin{enumerate}[resume]
\item A naive gateway counts an order as gone when it sends the cancel, and the strategy sends a replacement at once. What
position can the firm reach after one race, with a limit of 100 lots?
\item With $n$ races whose replacements are all still working, what is the worst over-position?
\item What does worst-case accounting do instead, and what does it cost the strategy?
\item At \qty{100}{\micro\second}, how many races does a strategy that cancels 40\,000 times a day expect?
\end{enumerate}

\medskip\noindent\textbf{Part IV --- The verdict.}
\begin{enumerate}[resume]
\item State the \emph{named result}: the race probability per cancel, and the worst over-position without in-flight
accounting.
\item Which two quantities does a firm control to reduce races, and which does it not?
\item Why does halving the one-way latency not halve the races when $d$ is large?
\item Why is a race not an error to be reported to the venue?
\item Why must the gateway keep a cancelled order in its table until the terminal report?
\item What does a replace change in the analysis?
\item How would you check, on a production day, that the races observed match the formula?
\item In one sentence: what is the rule that makes races harmless?
\end{enumerate}
\end{problem}

\section{Interview questions}

\begin{interviewq}[$\star$ \iqroles{developer}]\label{iq:ll:order-gateway-and-order-management:1}
What states does an order go through in an \omterm{def:ll:order-gateway-and-order-management:gateway}{order gateway}, and why are there more than at the exchange?
\end{interviewq}

\begin{interviewq}[$\star\star$ \iqroles{developer}]\label{iq:ll:order-gateway-and-order-management:2}
You send a cancel and receive a fill for the same order. What happened, and what should your system do?
\end{interviewq}

\begin{interviewq}[$\star\star$ \iqroles{developer}]\label{iq:ll:order-gateway-and-order-management:3}
Implement a rate limiter for an exchange session. What do you do with a cancel when the budget is exhausted?
\end{interviewq}

\begin{interviewq}[$\star\star$ \iqroles{developer, trader}]\label{iq:ll:order-gateway-and-order-management:4}
Your connection to the exchange drops for two seconds. What do you know about your orders, and what do you do when you
reconnect?
\end{interviewq}

\begin{interviewq}[$\star\star$ \iqroles{developer}]\label{iq:ll:order-gateway-and-order-management:5}
How do you check a position limit in constant time with thousands of open orders?
\end{interviewq}

\begin{interviewq}[$\star\star\star$ \iqroles{developer, researcher}]\label{iq:ll:order-gateway-and-order-management:6}
Design the order management layer for a market maker on three venues. What is the source of truth for positions, and how do
you know it is right?
\end{interviewq}
